\chapter[1938 --- Heymans]{1938: Corneille Heymans}
\label{ch:1938_corneilleheymans}

\section*{Main Contribution}

\textit{Discovered the role of the sinus and aortic body receptors in regulating respiration, explaining how the body detects changes in blood oxygen and carbon dioxide levels.}

\section{Official Citation}

for the discovery of the role of the sinus and aortic mechanisms in the regulation of respiration

\section{Background and Historical Context}

Corneille Jean François Heymans was born on March 28, 1892, in Ghent, Belgium, into a family with deep roots in the pharmaceutical and medical sciences. His father, Jean-François Heymans, was a professor of pharmacology at the University of Ghent and director of the university's pharmacological institute. Growing up in this intellectually stimulating environment, Corneille was exposed to scientific research from an early age, and he developed a keen interest in the mechanisms by which the human body regulates its vital functions.

Heymans pursued his medical education at the University of Ghent, where he excelled academically and demonstrated a particular aptitude for physiology and pharmacology. He graduated with his medical degree in 1914, just as World War I was engulfing Europe. During the war, Heymans served as a medical officer in the Belgian Army, an experience that exposed him to the devastating effects of trauma on the respiratory and cardiovascular systems. After the war, he returned to the University of Ghent, where he joined his father's laboratory and began his research career.

The scientific context of the 1930s was characterized by growing interest in the mechanisms of respiratory regulation. The work of several researchers had established that breathing is controlled by neural and chemical mechanisms, but the precise nature of these mechanisms remained poorly understood. The carotid bodies---small structures located at the bifurcation of the carotid arteries---were known to be involved in the regulation of breathing, but their exact function was unclear. It was against this backdrop that Heymans undertook his Nobel Prize-winning research.

\section{The Discovery/Contribution}

Corneille Heymans received the Nobel Prize in Physiology or Medicine ``for the discovery of the role of the sinus and aortic mechanisms in the regulation of respiration.'' His work provided the first definitive experimental demonstration that the carotid sinus and aortic arch contain specialized sensory receptors---now known as baroreceptors and chemoreceptors---that continuously monitor blood pressure and blood chemistry and relay this information to the brain to regulate breathing and cardiovascular function.

\subsection{The Carotid Sinus Reflex}

Heymans's primary contribution was the elucidation of the carotid sinus reflex, a fundamental homeostatic mechanism that maintains blood pressure within a narrow physiological range. The carotid sinus is a dilated region of the internal carotid artery located just above the bifurcation of the common carotid artery into its internal and external branches. Heymans demonstrated that this region contains stretch-sensitive nerve endings---baroreceptors---that detect changes in arterial blood pressure. When blood pressure rises, these receptors are stretched and send signals via the glossopharyngeal nerve (cranial nerve IX) to the cardiovascular and respiratory centers in the medulla oblongata of the brainstem. This input causes a reflex reduction in heart rate and blood vessel dilation, thereby lowering blood pressure back toward its normal level.

Conversely, when blood pressure falls, the baroreceptors are less stretched, reducing their firing rate. This decreased input to the brainstem leads to increased heart rate and vasoconstriction, thereby raising blood pressure. Heymans's meticulous experimental work, which involved the ingenious technique of transplanting the carotid sinus and its nerve connections to different locations in the body and demonstrating that the reflex still functioned, provided compelling evidence for the existence and importance of this feedback mechanism.

\subsection{The Aortic Arch and Respiratory Regulation}

In addition to the carotid sinus, Heymans demonstrated that the aortic arch contains similar sensory receptors that contribute to the regulation of respiration and cardiovascular function. The aortic bodies---small clusters of chemosensitive cells located in the wall of the aortic arch---detect changes in the chemical composition of the blood, particularly the oxygen and carbon dioxide levels. When oxygen levels in the blood fall (hypoxia) or carbon dioxide levels rise (hypercapnia), these chemoreceptors send signals via the vagus nerve (cranial nerve X) to the brainstem respiratory center, stimulating an increase in the rate and depth of breathing.

Heymans's work on the aortic chemoreceptors was particularly significant because it demonstrated that the body has multiple, redundant mechanisms for monitoring and regulating blood gas levels. This redundancy ensures that the vital process of respiration is maintained even if one set of receptors is compromised. The aortic bodies, while less sensitive than the carotid bodies, provide an additional source of information about blood chemistry that helps fine-tune the respiratory response.

\subsection{Experimental Innovation}

One of a notable aspects of Heymans's research was his experimental methodology. To demonstrate that the carotid sinus and aortic arch were the sites of the sensory receptors responsible for respiratory and cardiovascular regulation, he developed a technique that has since become known as ``cross-circulation'' experiments. In these experiments, Heymans would transplant the carotid sinus of one animal into the circulatory system of another animal, carefully connecting the blood vessels and nerves so that the transplanted sinus received blood from the host animal while its nerve connections to the brainstem remained intact. This elegant approach allowed him to isolate the effects of the sensory receptors from other physiological variables and provided definitive proof that the carotid sinus and aortic arch were the primary sites of blood pressure and blood gas sensing.

Heymans also conducted experiments in which he selectively stimulated or denervated the carotid sinus and aortic bodies, measuring the resulting changes in breathing and cardiovascular function. These experiments allowed him to construct a detailed picture of the neural pathways involved in respiratory regulation and to establish the principles of negative feedback that govern this vital homeostatic system.

\section{Impact on Modern Medicine}

Heymans's discoveries have had a significant impact on modern medicine, particularly in the fields of cardiology, anesthesiology, and critical care medicine. The understanding of baroreceptor reflexes that emerged from his work is fundamental to the clinical management of hypertension (high blood pressure). Modern antihypertensive medications are designed to work with, rather than against, these natural regulatory mechanisms. For example, beta-blockers, which are among the most widely prescribed drugs for hypertension and heart failure, act in part by modifying the heart's response to sympathetic nervous system input---a system that is regulated, in part, by the baroreceptor reflex that Heymans elucidated.

The concept of chemoreceptor regulation of breathing has become central to the management of patients with respiratory failure. In intensive care units around the world, patients on mechanical ventilators are monitored for changes in blood oxygen and carbon dioxide levels that reflect the function of the chemoreceptors Heymans described. Understanding how these receptors respond to changes in blood gases has also been crucial for the development of supplemental oxygen therapy and for the management of patients with chronic obstructive pulmonary disease (COPD) and other chronic respiratory conditions.

Heymans's work also laid the groundwork for our understanding of syncope (fainting), a common clinical condition in which a sudden drop in blood pressure leads to temporary loss of consciousness. The carotid sinus reflex, when overly sensitive, can cause a sudden and dramatic drop in heart rate and blood pressure---a condition known as carotid sinus hypersensitivity---that can trigger fainting episodes. Recognition and diagnosis of this condition are based directly on the physiological principles that Heymans established.

In the field of anesthesiology, knowledge of the baroreceptor reflex is essential for the safe management of patients undergoing surgery. General anesthesia can suppress the baroreceptor reflex, leading to dangerous fluctuations in blood pressure during and after surgical procedures. Anesthesiologists must be intimately familiar with the mechanisms Heymans described in order to anticipate and manage these changes, ensuring patient safety throughout the perioperative period.

More recently, Heymans's discoveries have informed research into autonomic dysfunction in conditions such as diabetes mellitus, heart failure, and Parkinson's disease. In these conditions, the baroreceptor reflex may be impaired, leading to orthostatic hypotension---a sudden drop in blood pressure upon standing that can cause dizziness, falls, and syncope. Understanding the pathophysiology of this condition, which is essential for developing effective treatments, builds directly on the foundation that Heymans laid in the 1930s.

The principles of negative feedback that Heymans described in the context of respiratory and cardiovascular regulation have also had broader implications for systems biology and control theory. The baroreceptor reflex is one of the best-understood examples of a negative feedback system in biology, and it has served as a model for understanding other homeostatic mechanisms, including thermoregulation, glucose homeostasis, and hormone regulation. The conceptual framework that Heymans established---in which sensory receptors detect deviations from a set point and trigger corrective responses---has become a cornerstone of modern physiology.

\section{Legacy}

Corneille Heymans's legacy extends beyond his specific discoveries to encompass his broader contributions to physiology and medicine. He served as professor of pharmacology at the University of Ghent from 1930 until his retirement, and during this long career, he trained dozens of students who went on to distinguished careers in the biomedical sciences. His laboratory at the University of Ghent became a world-renowned center for the study of cardiovascular and respiratory physiology, attracting researchers from around the globe.

Heymans was also a tireless advocate for international scientific cooperation. He served as president of the International Union of Physiological Sciences and was a member of numerous scientific organizations around the world. He received honorary degrees from universities in many countries and was elected to prestigious scientific academies, including the Royal Society of London and the National Academy of Sciences of the United States.

The principles that Heymans elucidated continue to be taught in medical schools and physiology courses around the world. The baroreceptor reflex and chemoreceptor regulation of breathing are fundamental topics in every medical curriculum, and the experimental techniques that Heymans pioneered---particularly the cross-circulation technique---remain important tools in cardiovascular and respiratory research. As our understanding of autonomic physiology continues to deepen, the foundational work of Corneille Heymans remains an essential reference point and an enduring source of inspiration for scientists and clinicians alike.


\section{Modern Developments}

Heymans' work on carotid body physiology has informed modern respiratory medicine. Understanding of chemoreceptor function has led to treatments for sleep apnea (CPAP) and chronic obstructive pulmonary disease. Carotid body tumors are now diagnosed and treated with precision surgery. Understanding of oxygen sensing has revealed the HIF (hypoxia-inducible factor) pathway, which won the 2019 Nobel Prize and has led to treatments for anemia (roxadustat, daprodustat) and is being explored for cancer therapy.
